\section{Pieces with a small boundary}\label{sec:affine}

This section proves that two disjoint pieces of a cubic graph, each large, with a small boundary and
a large girth, force $q(G,P)\le\frac9{25}$ (Proposition~\ref{prop:blocks}). The only input is
Theorem~\ref{thm:forest}, in the following form.

\begin{proposition}[Forest input]\label{prop:input}
There is an integer $R_{*}\ge2^{100}$ with the following property. Let $H$ be a connected multigraph
of rank $\rho\ge R_{*}$, and let $W$ be a cycle of $H$ with at least three and at most
$\rho/\log^{2}\rho$ edges. Then
\[
  \PP\bigl(\xi_H\sm E(W)\text{ is a forest}\bigr)\;\le\;\frac{11}{20}.
\]
\end{proposition}

\begin{proof}
Replace every edge of $H$ outside $W$ by a path with three edges. The new graph $H'$ is simple and
connected, it contains $W$, and it has rank $\rho$. The even subgraphs of $H'$ correspond to those
of $H$, and one has a forest outside $W$ if and only if the other has. Apply
Theorem~\ref{thm:forest} to $H'$ and $W$, in the range of Remark~\ref{rem:range}: $W$ is one cycle
with at most $\rho/\log^{2}\rho$ edges. Choose $R_{*}$ so large that $\rho\ge R_{*}$ is in the range
of the theorem, $\log\rho\ge100$, and $20/\lgk3\rho\le\frac1{20}$.
\end{proof}

\paragraph{Status of the input.}
Theorem~\ref{thm:forest} is formalised for $|E(W)|\le\sqrt r$. The range $|E(W)|\le r/\log^{2}r$
is Remark~\ref{rem:range}, which is not formalised. Lemma~\ref{lem:width} below needs the larger
range.

\subsection{Localisation}

\begin{lemma}[Localisation]\label{lem:loc}
Let $G$ be a connected multigraph, and let $X$ be a nonempty proper set of vertices such that
$G-X$ is connected. Put $d=|\delta X|$. Form $G_X$ from $G[X]$ as follows. For each edge of
$\delta X$ add a new vertex and join it to the end of that edge in $X$. Join the new vertices by a
cycle $W_X$ through all of them; if $d<3$, first add $3-d$ further new vertices on that cycle. Then
\begin{enumerate}
\item[\textup{(a)}] $G_X$ is connected, and $\beta(G_X)=\beta(G[X])+d-\comp(G[X])+1$;
\item[\textup{(b)}] if every vertex of $X$ has degree three in $G$, then
$\beta(G_X)=(|X|+d)/2+1$;
\item[\textup{(c)}] the set $\xi\cap E(G[X])$ has the same distribution for $\xi$ uniform on
$\Z(G)$ and for $\xi$ uniform on $\Z(G_X)$;
\item[\textup{(d)}] an even subgraph $\xi$ of $G_X$ has a forest outside $E(W_X)$ if and only if
$\xi\cap E(G[X])$ is a forest.
\end{enumerate}
\end{lemma}

\begin{proof}
(a) Every component of $G[X]$ has an edge to $G-X$, because $G$ is connected. So the cycle $W_X$
joins all components. Without the extra vertices, $G_X$ has $d$ more vertices and $2d$ more edges
than $G[X]$, and it is connected. Extra vertices on the cycle add as many edges as vertices.

(b) Count the edge ends at $X$: $3|X|=2|E(G[X])|+d$. With (a) this gives the formula.

(c) Contract the connected set $V(G)\sm X$ of $G$ to one vertex, and delete the loops that arise.
Call the result $G/X$. An even subgraph of $G$ restricts to an even subgraph of $G/X$. Every even
subgraph of $G/X$ arises in this way: the missing edges inside the contracted set can be chosen in
a spanning tree of that set. So the restriction is a surjective linear map $\Z(G)\to\Z(G/X)$, and it
sends the uniform distribution to the uniform distribution. Now contract $W_X$ in $G_X$. The result
is the same multigraph $G/X$, with the new edges in the place of the edges of $\delta X$. So both
distributions are the image of the uniform distribution on $\Z(G/X)$.

(d) The edges of $G_X$ outside $W_X$ are the edges of $G[X]$ and the $d$ new edges at $X$. In
$G_X-E(W_X)$ each of these new edges has an end of degree one. So it lies on no cycle of
$G_X-E(W_X)$.
\end{proof}

A version of this lemma for connected $X$ was the starting point. The version above, for a set $X$
with a connected complement, is due to Astra.

\subsection{Thin markers force thin layers}

\begin{lemma}[Layers]\label{lem:width}
Put $B=2R_{*}+2^{100}+1$. Let $G$ be a connected subcubic multigraph, and let $P$ be a nonempty set
of vertices such that $G[P]$ is connected and every vertex outside $P$ has degree three. Fix a
vertex $o$, and call the set of vertices at distance $j$ from $o$ the \emph{layer} $j$. Suppose that
$P$ has at most $k\ge1$ vertices in each layer, and that $q(G,P)\ge\frac35$. Then every layer has at
most $B(k+1)^{2}$ vertices.
\end{lemma}

\begin{proof}
Put $\Phi(m)=m/\log^{2}m$ and $\Theta=\max\{2R_{*},\,2^{100}(k+1)^{2}\}$.

\emph{Step 1: no large set outside $P$ has a small boundary.} We claim that no set
$X\subseteq V(G)\sm P$ has $|X|=m\ge\Theta$ and $|\delta X|\le\Phi(m)/4$. Suppose that $X$ is such a
set. Let $C_0$ be the component of $G-X$ that contains $P$, and put $X^{+}=V(G)\sm V(C_0)$. Then
$X\subseteq X^{+}$, the set $X^{+}$ avoids $P$, and $G-X^{+}=C_0$ is connected. Different components
of $G-X$ have no edge between them, so $\delta X^{+}\subseteq\delta X$. Every vertex of $X^{+}$ has
degree three. Let $d^{+}=|\delta X^{+}|\ge1$ and apply Lemma~\ref{lem:loc} to $X^{+}$. The rank of
$G_{X^{+}}$ is
\[
  \rho=\frac{|X^{+}|+d^{+}}2+1\;\ge\;\frac m2\;\ge\;R_{*}.
\]
The cycle $W_{X^{+}}$ has $\max\{d^{+},3\}$ edges. The function $\Phi$ is increasing on
$[8,\infty)$, and $\Phi(m)/4\le\Phi(m/2)$. So
$d^{+}\le\Phi(m)/4\le\Phi(m/2)\le\Phi(\rho)$, and also $3\le\Phi(\rho)$. By
Proposition~\ref{prop:input} and
Lemma~\ref{lem:loc}(c),(d),
\[
  \PP\bigl(\xi_G\cap E(G[X^{+}])\text{ is a forest}\bigr)\;\le\;\frac{11}{20}.
\]
But $G[X^{+}]$ is a subgraph of $G-P$, so this probability is at least $q(G,P)\ge\frac35$. This
contradiction proves the claim.

\emph{Step 2: a window of layers.} Let $m_j$ be the number of vertices of layer $j$ outside $P$, with
$m_j=0$ outside the range of the layers. For $i\le i'$ let $X(i,i')$ be the set of vertices outside
$P$ in the layers $i,\dots,i'$. An edge joins vertices of the same layer or of consecutive layers.
An edge that leaves $X(i,i')$ ends outside $P$ in layer $i-1$ or $i'+1$, or it ends in $P$ in one of
the layers $i-1,\dots,i'+1$. So
\begin{equation}\label{eq:window}
  |\delta X(i,i')|\;\le\;3m_{i-1}+3m_{i'+1}+3k(i'-i+3).
\end{equation}
Suppose that $m_j\ge\Theta$ for some $j$. Start with the window $i=i'=j$. As long as a layer next
to the window has more than $m/(48\log^{2}m)$ vertices outside $P$, where $m=|X(i,i')|$ is the
present size, add that layer to the window. The graph is finite, so this stops. Let $\mu$ be the
final size and $\lambda=\log\mu\ge100$. Each step multiplies the size by more than
$1+1/(48\lambda^{2})$. So the number of steps is at most
\[
  \frac{\ln\mu}{\ln\bigl(1+1/(48\lambda^{2})\bigr)}\;\le\;96\lambda^{2}\ln\mu\;<\;67\lambda^{3}.
\]
At the end both layers next to the window have at most $\mu/(48\lambda^{2})$ vertices outside $P$.
By~\eqref{eq:window},
\[
  |\delta X(i,i')|\;\le\;\frac{\mu}{8\lambda^{2}}+3k\bigl(67\lambda^{3}+3\bigr)
  \;\le\;\frac{\mu}{4\lambda^{2}}=\frac{\Phi(\mu)}4 .
\]
The second inequality says $24k\lambda^{2}(67\lambda^{3}+3)\le2^{\lambda}$. It holds because
$k\le2^{\lambda/2-50}$, by $\mu\ge\Theta\ge2^{100}(k+1)^{2}$, and because
$24\lambda^{2}(67\lambda^{3}+3)\le2000\lambda^{5}\le2^{\lambda/2+50}$ for $\lambda\ge100$. So the
final window contradicts Step~1.

Hence every layer has fewer than $\Theta$ vertices outside $P$ and at most $k$ vertices in $P$. Their
number is at most $\Theta+k\le B(k+1)^{2}$.
\end{proof}

\subsection{Poles and their forest proportions}

A \emph{pole} is a connected multigraph in which every vertex has degree two or three. Its
vertices of degree two are its \emph{ports}. Let $N$ be a pole and let $T$ be a set of ports of even
size. A \emph{$T$-join} of $N$ is a set of edges in which the vertices of $T$ have odd degree and
all other vertices have even degree. There are $2^{\beta(N)}$ of them. The \emph{forest proportion}
of $T$ is the proportion of the $T$-joins that are forests. For $T=\emptyset$ the $T$-joins are the
even subgraphs, only the empty one is a forest, and the forest proportion is $2^{-\beta(N)}$.

\begin{lemma}[A marker in a pole]\label{lem:marker}
Let $N$ be a pole with $k\ge2$ ports. Suppose that some set $T$ of ports has forest proportion at
least $\frac35$. Then, for every port $o$, each layer of $N$ with respect to $o$ has at most
$B(k+1)^{2}$ vertices.
\end{lemma}

\begin{proof}
Take a breadth-first tree of $N$ from $o$, and let $P$ be the union of its paths from $o$ to the
other ports. Then $N[P]$ is connected, $P$ contains all vertices of degree two, and $P$ has at most
$k$ vertices in each layer, because a path of the tree from $o$ has one vertex in each layer that
it meets. The tree paths inside $P$ contain a $T$-join $F_0$. The map $J\mapsto J+F_0$ is a
bijection from the $T$-joins of $N$ to the even subgraphs of $N$, and $J+F_0$ has the same edges in
$N-P$ as $J$. So if $J$ is a forest, then $J+F_0$ is a forest in $N-P$. Hence
$q(N,P)\ge\frac35$, and Lemma~\ref{lem:width} applies to $N$ and $P$.
\end{proof}

\begin{lemma}[Girth and rank of a pole]\label{lem:moore}
Let $N$ be a pole of rank $h\ge2$ with $k$ ports and girth at least $D$. Put
$t=\lceil\log(8h)\rceil$. Then
\[
  D\;\le\;4t\Bigl(1+\frac{k}{3(h-1)}\Bigr).
\]
\end{lemma}

\begin{proof}
Replace every vertex of degree two and its two edges by one edge, and give each edge of the result
the length of the path of $N$ that it replaces. The result $N_0$ is a connected cubic multigraph
with $n_0=2h-2$ vertices and $m_0=3h-3$ edges, and the lengths of its edges add up to $m_0+k$. A
cycle of $N_0$ comes from a cycle of $N$ of the same length.

There are $2m_0\,2^{t-1}$ non-backtracking walks with $t$ edges in $N_0$. Choose one of them
uniformly. Then each of its steps is uniform over the $2m_0$ directed edges. So the expected length
of the walk is $t(1+k/m_0)$. At least half of the walks have length at most twice this value. Their
number is at least
\[
  m_0\,2^{t-1}\;\ge\;12h(h-1)\;>\;n_0^{2}.
\]
If one of these short walks visits a vertex twice, it contains a cycle that is not longer than the
walk. Otherwise all of them are paths, and two of them have the same first vertex and the same last
vertex. Their union contains a cycle of length at most the sum of their lengths. In both cases $N$
has a cycle of length at most $4t(1+k/m_0)$.
\end{proof}

\begin{lemma}[Forest proportions at large girth]\label{lem:affine}
There are absolute constants $D_{*}\ge3$ and $C_0$ with the following property. Let $N$ be a pole
of rank $h$ with $k$ ports and girth at least $D$. Suppose that
\begin{enumerate}
\item[\textup{(a)}] $h\ge256$, $k\le4D$ and $D\ge D_{*}$, or
\item[\textup{(b)}] $h\ge\Lambda\ge4$, $k\le2\Lambda$ and $D\ge C_0\log(\Lambda+2)$ for some
number $\Lambda$.
\end{enumerate}
Then every set of ports of even size has forest proportion at most $\frac35$.
\end{lemma}

\begin{proof}
If $k\le1$, the only set is the empty one, with forest proportion $2^{-h}$. Let $k\ge2$, and put
$t=\lceil\log(8h)\rceil\le\log h+4$.

\emph{The rank is exponential in the girth.} In case~(a), $h-1>\frac{32}3t$: this holds at
$h=256$, where $t=11$, and the ratio $(h-1)/(\log h+4)$ increases with $h$. So
Lemma~\ref{lem:moore} gives $D\le4t+16tD/(3(h-1))<4t+D/2$. In case~(b),
$k/(3(h-1))\le2h/(3(h-1))<1$, and Lemma~\ref{lem:moore} gives $D<8t$. In both cases
$D<8t\le8\log h+32$, so
\begin{equation}\label{eq:rankgirth}
  h\;\ge\;2^{D/8-4},\qquad\text{and }N\text{ has }n\ge h\text{ vertices}.
\end{equation}

\emph{A ball without ports.} Put $s=\lfloor D/16\rfloor$. A ball of radius $s$ has at most
$3\cdot2^{s}$ vertices. So at most $3k\,2^{s}$ vertices are within distance $s$ of a port. In
case~(a) this is at most $12D\,2^{D/16}$, and in case~(b) at most $6\Lambda\,2^{D/16}$. Choose
$D_{*}$ and $C_0$ so large that this number is less than $2^{D/8-4}\le n$ in each case. Then some
vertex $w$ has distance more than $s$ from every port. All vertices within distance $s$ of $w$
have degree three, and $2s<D$. So the ball of radius $s$ around $w$ is a tree in which every
vertex at distance less than $s$ from $w$ has degree three, and it has at least
$2^{s}\ge2^{D/16-1}$ vertices. Fix a port $o$. The distances from $o$ to two vertices of this ball
differ by at most $2s$. So the ball meets at most $2s+1\le D$ layers with respect to $o$, and some
layer has at least $2^{D/16-1}/D$ vertices.

\emph{Conclusion.} Suppose that some set of ports has forest proportion more than $\frac35$. By
Lemma~\ref{lem:marker} every layer has at most $B(k+1)^{2}$ vertices. This is at most
$B(4D+1)^{2}$ in case~(a) and at most $B(2\Lambda+1)^{2}$ in case~(b). In case~(a) choose $D_{*}$ so
large that $2^{D/16-1}/D>B(4D+1)^{2}$ for $D\ge D_{*}$. In case~(b) the function $2^{D/16}/D$
increases for $D\ge32$, and at $D=C_0\log(\Lambda+2)$ its ratio to $(\Lambda+2)^{2}$ is
$(\Lambda+2)^{C_0/16-2}/(C_0\log(\Lambda+2))$. This tends to infinity with $C_0$, uniformly in
$\Lambda\ge4$. So a large $C_0$ gives $2^{D/16-1}/D>B(2\Lambda+1)^{2}$. In both cases we have a
contradiction.
\end{proof}

\subsection{Two pieces}

\begin{lemma}[Product]\label{lem:product}
Let $G$ be a cubic multigraph, let $P$ be a set of vertices, and let $U_1,U_2\subseteq V(G)\sm P$
be disjoint sets such that $G[U_1]$ and $G[U_2]$ are poles. Suppose that in each of these poles
every set of ports of even size has forest proportion at most $\frac35$. Then
$q(G,P)\le\frac9{25}$.
\end{lemma}

\begin{proof}
Let $E'$ be the set of the edges of $G$ that lie neither in $G[U_1]$ nor in $G[U_2]$. Fix a set
$x\subseteq E'$ that is the restriction of some even subgraph. For $i=1,2$ let $T_i$ be the set of
the vertices of $U_i$ that meet an odd number of edges of $x$. An edge set $\xi$ with
$\xi\cap E'=x$ is even if and only if $\xi\cap E(G[U_i])$ is a $T_i$-join of $G[U_i]$ for $i=1,2$:
the degree condition at a vertex of $U_i$ involves only edges of $G[U_i]$ and of $x$, and at the
other vertices it involves only $x$. So, given $\xi\cap E'=x$, the two sets $\xi\cap E(G[U_i])$ are
independent, and each is a uniform $T_i$-join. A vertex of degree three in $G[U_i]$ meets no edge
of $E'$, because $G$ is cubic. So $T_i$ is a set of ports, of even size.

If $\xi\cap E(G-P)$ is a forest, then $\xi\cap E(G[U_1])$ and $\xi\cap E(G[U_2])$ are forests. Given
$\xi\cap E'=x$, this has probability at most $\frac35\cdot\frac35$. Average over $x$.
\end{proof}

Edges between $U_1$ and $U_2$ belong to $E'$. So the two sets may be adjacent.

\begin{lemma}[The core of a piece]\label{lem:core}
Let $G$ be a cubic multigraph, and let $X$ be a set of vertices such that $G[X]$ is connected and
has rank $h\ge2$. Then there is a set $U\subseteq X$ such that $G[U]$ is a pole of rank $h$ with at
most $|\delta X|$ ports, and the girth of $G[U]$ is at least the girth of $G[X]$.
\end{lemma}

\begin{proof}
Delete vertices of degree one from $G[X]$ as long as there are any. The result $G[U]$ is connected,
has the same cycles as $G[X]$ and all its degrees are two or three. When a vertex of degree one is
deleted, two edges that left the set no longer do, and one edge inside the set now leaves it. So
$|\delta U|\le|\delta X|$. A vertex of degree two in $G[U]$ has exactly one edge end outside
$G[U]$, and a vertex of degree three has none. So $G[U]$ has $|\delta U|$ ports.
\end{proof}

\begin{proposition}[Two pieces]\label{prop:blocks}
Let $G$ be a connected cubic multigraph and let $P$ be a set of vertices. Let
$X_1,X_2\subseteq V(G)\sm P$ be disjoint, nonempty and proper sets of vertices. Suppose that one of
the following holds for $i=1,2$.
\begin{enumerate}
\item[\textup{(a)}] $G[X_i]$ is connected of rank at least $256$ and girth at least $D\ge D_{*}$,
and $|\delta X_i|\le4D$.
\item[\textup{(b)}] $|X_i|\ge8\Lambda^{2}$, $|\delta X_i|\le2\Lambda$, and $G[X_i]$ has girth at
least $D$, where $\Lambda\ge4$ and $D\ge C_0\log(\Lambda+2)$.
\end{enumerate}
Then $q(G,P)\le\frac9{25}$.
\end{proposition}

\begin{proof}
It suffices to find sets $U_i\subseteq X_i$ as in Lemma~\ref{lem:product}. In case~(a),
Lemma~\ref{lem:core} gives a pole $G[U_i]$ of rank at least $256$ with at most $4D$ ports and girth
at least $D$, and Lemma~\ref{lem:affine}(a) applies.

In case~(b) write $n=|X_i|$, $d=|\delta X_i|$ and $c=\comp(G[X_i])$. Every component of $G[X_i]$
has an edge that leaves $X_i$, so $1\le c\le d\le2\Lambda$. Count the edge ends:
$3n=2|E(G[X_i])|+d$. So $\beta(G[X_i])=(n-d)/2+c$, and some component $Y$ of $G[X_i]$ has rank at
least
\[
  \frac{n-d}{2c}+1\;\ge\;\frac{8\Lambda^{2}-2\Lambda}{4\Lambda}+1\;>\;\Lambda .
\]
Its boundary is a part of $\delta X_i$. Lemma~\ref{lem:core} for $Y$ gives a pole $G[U_i]$ of rank
at least $\Lambda$ with at most $2\Lambda$ ports and girth at least $D$, and
Lemma~\ref{lem:affine}(b) applies.
\end{proof}
